\documentclass[11pt,a4paper]{article}
\usepackage[utf8]{inputenc}
\usepackage[margin=1in]{geometry}
\usepackage{enumitem}
\usepackage{booktabs}
\usepackage{hyperref}

\title{\textbf{Course Breakdown: AS.110.106 Calculus I (Bio. \& Soc. Sci.)}}
\author{\textbf{Department of Mathematics, Johns Hopkins University}}
\date{}

\begin{document}

\maketitle

\section*{Course Overview}
\begin{itemize}
    \item \textbf{Course Title:} Calculus I (Biology and Social Sciences)[cite: 1]
    \item \textbf{Course Code:} AS.110.106[cite: 1]
    \item \textbf{Textbook:} \textit{Calculus for Biology and Medicine}, 4th Edition, C. Neuhauser and M. Roper (Pearson, 2018)[cite: 1]
\end{itemize}

---

\section*{1. Prerequisite Topics}
To succeed in this course, students should be familiar with:
\begin{itemize}
    \item \textbf{Algebra:} High school algebra, solving equations, exponential/logarithmic rules, and polynomial factoring.
    \item \textbf{Trigonometry:} Basic trigonometric functions ($\sin, \cos, \tan$), the unit circle, and standard trigonometric identities.
    \item \textbf{Basic Precalculus:} Graphing elementary functions, domain and range analysis, and understanding functional notation.
\end{itemize}

---

\section*{2. Course Structure \& Timeline}
The course is structured into 5 core modules delivered across approximately 12+ weeks of instruction[cite: 1]:

\begin{center}
\begin{tabular}{@{}lll@{}}
\toprule
\textbf{Module} & \textbf{Core Focus} & \textbf{Estimated Duration} \\ \midrule
Module 1 & Review of Basic Properties of Functions \& Sequences & 2 Weeks[cite: 1] \\
Module 2 & Limits and Continuity & 2 Weeks[cite: 1] \\
Module 3 & Derivatives \& Differentiation Rules & 3 Weeks[cite: 1] \\
Module 4 & Applications of Differentiation & 2+ Weeks[cite: 1] \\
Module 5 & Integration \& Integration Techniques & 3+ Weeks[cite: 1] \\ \bottomrule
\end{tabular}
\end{center}

---

\section*{3. Detailed Modules, Chapters, Concepts, \& Objectives}

\subsection*{Module 1: Basic Properties of Functions and Sequences (2 Weeks) \cite{1}}
\begin{itemize}
    \item \textbf{Chapters / Sections Covered:} Chapter 1, Section 2.1, Section 2.2[cite: 1]
    \item \textbf{Key Concepts:}
    \begin{itemize}
        \item Linear, polynomial, exponential, and logarithmic functions in biological/social contexts.
        \item Exponential growth and decay models[cite: 1].
        \item Discrete-time models, sequences, and recursion relations[cite: 1].
    \end{itemize}
    \item \textbf{Learning Objectives:}
    \begin{itemize}
        \item Model population growth, radioactive decay, and resource consumption using exponential functions.
        \item Evaluate terms and limits of discrete sequences[cite: 1].
    \end{itemize}
\end{itemize}

\subsection*{Module 2: Limits and Continuity (2 Weeks) \cite{1}}
\begin{itemize}
    \item \textbf{Chapters / Sections Covered:} Sections 3.1, 3.2, 3.3, 3.4, 3.5, 3.6[cite: 1]
    \item \textbf{Key Concepts:}
    \begin{itemize}
        \item Intuitive and formal definitions ($\epsilon$-$\delta$) of limits (Sections 3.1, 3.6)[cite: 1].
        \item Continuity and properties of continuous functions (Sections 3.2, 3.5)[cite: 1].
        \item Asymptotes and limits at infinity (Section 3.3)[cite: 1].
        \item Trigonometric limits and the Sandwich (Squeeze) Theorem (Section 3.4)[cite: 1].
    \end{itemize}
    \item \textbf{Learning Objectives:}
    \begin{itemize}
        \item Calculate limits algebraically and graphically.
        \item Apply the Sandwich Theorem to determine limits of bounded oscillating functions[cite: 1].
        \item Identify types of discontinuities and apply properties of continuous functions (e.g., Intermediate Value Theorem)[cite: 1].
    \end{itemize}
\end{itemize}

\subsection*{Module 3: Derivatives (3 Weeks) \cite{1}}
\begin{itemize}
    \item \textbf{Chapters / Sections Covered:} Sections 4.1 -- 4.11[cite: 1]
    \item \textbf{Key Concepts:}
    \begin{itemize}
        \item Formal definition of the derivative as a limit of difference quotients (Section 4.1)[cite: 1].
        \item Power, constant, sum, product, quotient, and chain rules (Sections 4.2--4.5)[cite: 1].
        \item Derivatives of trigonometric, exponential, and inverse functions (Sections 4.8--4.10)[cite: 1].
        \item Implicit differentiation and higher-order derivatives (Sections 4.6, 4.7)[cite: 1].
        \item Linear approximations and differentials (Section 4.11)[cite: 1].
    \end{itemize}
    \item \textbf{Learning Objectives:}
    \begin{itemize}
        \item Differentiate complex explicit and implicit equations.
        \item Interpret the derivative as an instantaneous rate of change in biological systems.
        \item Use tangent line approximations to estimate function values[cite: 1].
    \end{itemize}
\end{itemize}

\subsection*{Module 4: Applications of Differentiation (2+ Weeks) \cite{1}}
\begin{itemize}
    \item \textbf{Chapters / Sections Covered:} Sections 5.1, 5.2, 5.3, 5.4, 5.5[cite: 1]
    \item \textbf{Key Concepts:}
    \begin{itemize}
        \item Extreme Value Theorem and Mean Value Theorem (Section 5.1)[cite: 1].
        \item First and second derivative tests for monotonicity, concavity, and inflection points (Sections 5.2, 5.3)[cite: 1].
        \item Optimization problems in life and social sciences (Section 5.4)[cite: 1].
        \item L'Hôpital's Rule for indeterminate forms (Section 5.5)[cite: 1].
    \end{itemize}
    \item \textbf{Learning Objectives:}
    \begin{itemize}
        \item Sketch accurate curves using first and second derivatives[cite: 1].
        \item Solve real-world optimization problems (e.g., maximizing yield or minimizing cost)[cite: 1].
        \item Resolve indeterminate limits using L'Hôpital's Rule[cite: 1].
    \end{itemize}
\end{itemize}

\subsection*{Module 5: Integration and Techniques of Integration (3+ Weeks) \cite{1}}
\begin{itemize}
    \item \textbf{Chapters / Sections Covered:} Sections 5.10, 6.1, 6.2, 6.3, 7.1, 7.2, 7.3[cite: 1]
    \item \textbf{Key Concepts:}
    \begin{itemize}
        \item Antiderivatives and net change (Section 5.10)[cite: 1].
        \item Riemann sums and the Definite Integral (Section 6.1)[cite: 1].
        \item The Fundamental Theorem of Calculus (Section 6.2)[cite: 1].
        \item Substitution Rule and Integration by Parts (Sections 7.1, 7.2)[cite: 1].
        \item Partial fractions for rational functions (Section 7.3)[cite: 1].
    \end{itemize}
    \item \textbf{Learning Objectives:}
    \begin{itemize}
        \item Compute exact areas under curves using antiderivatives[cite: 1].
        \item Select and apply appropriate techniques of integration (Substitution, Integration by Parts, Partial Fractions) to solve integral problems[cite: 1].
        \item Apply integration to physical, biological, and social models[cite: 1].
    \end{itemize}
\end{itemize}

---

\section*{4. Difficult Topics Identified}
Students often struggle with the following concepts in this syllabus:
\begin{enumerate}
    \item \textbf{Formal $\epsilon$-$\delta$ Definitions (Section 3.6):} Moving from intuitive graphical limit concepts to rigorous mathematical proofs[cite: 1].
    \item \textbf{Implicit Differentiation \& Related Rates (Section 4.6):} Applying the chain rule to multi-variable relations where variables are implicitly dependent on time or another variable[cite: 1].
    \item \textbf{Word Problems in Optimization (Section 5.4):} Translating descriptive word problems from biological or social contexts into symbolic objective equations[cite: 1].
    \item \textbf{Advanced Integration Techniques (Sections 7.1--7.3):} Identifying whether to use Substitution, Integration by Parts, or Partial Fractions for a given integral[cite: 1].
\end{enumerate}

\end{document}